% Source: https://artofproblemsolving.com/community/c3h1911042_mathematical_puzzle
% Stable identity: aops-topic-1911042-q01
\textbf{$n=5$ works.} Let the numbers be $a_1,\dots,a_5$ and consider the five sums $s_k=a_1+\cdots+a_k$ ($k=1,\dots,5$). If some $s_k$ is a multiple of $5$, we are done. Otherwise the five sums have remainders among $1,2,3,4$, so by the pigeonhole principle two of them, say $s_i$ and $s_j$ with $i<j$, have the same remainder. Then $s_j-s_i=a_{i+1}+\cdots+a_j$ is a multiple of $5$.

\textbf{$n=4$ does not.} For $1,6,11,16$, every number leaves remainder $1$ on division by $5$, so a sum of $k$ of them ($1\le k\le4$) leaves remainder $k$, never $0$.

So the answer is $\boxed{5}$.
